\section{Approach}
\label{sec:mrmr-approach}
Super Odometry is modeled using a factor graph. Here we mainly introduce three types of odometry factors for factor graph construction, namely: (a) IMU; (b) LiDAR-inertial; and (c) Visual-inertial odometry factors. 

\subsection{IMU Odometry Factors}
The structure of IMU odometry factors is shown in Fig.\ref{fig:algorithm}.({a}), which is similar to traditional pose graph structure. However, it is different since the estimated $state$ of IMU odometry not only contains position $\mathbf{p}_{b}^{w}$, velocity $\mathbf{v}_{b}^{w}$ and orientation 
$\mathbf{q}_{b}^{w}$ but also the accelerometer $\mathbf{b}_{a}$ and gyroscope bias $\mathbf{b}_{g}$. Each node in the pose graph is associated with a state. 
The density of nodes is determined by the lowest-frequency odometry. The edge between two consecutive nodes represents the relative body motion obtained from 
IMU preintegration method. Other edges could be local constraints or global constraints which depend on the type of sensors. 

\subsubsection{\textbf{IMU Preintegration Factor}}
Based on the IMU preintegration method proposed in \cite{forster2015imu}, we can obtain the relative motion measurement $\Delta \mathbf{r}=\left(\hat{\alpha}_{j}^{i}, \hat{\beta}_{j}^{i}, \hat{\gamma}_{j}^{{i}}\right)$ between consecutive IMU body frames $i$ and $j$. The IMU preintegration factor can be defined as:
\begin{equation}
	\label{eq:imu_factor}
	\begin{aligned}
		\mathbf{e}_{ij}^{{imu}}&=\mathbf{z}_{ij}^{b}-h_{ij}^{b}\left(\mathbf{x}_{i}, \mathbf{x}_{j}\right) \\
		&=\left[\begin{array}{c}
			\hat{\alpha}_{j}^{i}\\
			\hat{\beta}_{j}^{i} \\
			\hat{\gamma}_{j}^{i} \\
			\mathbf{b}_{a {j}} \\
			\mathbf{b}_{w {j}}\\ 
		\end{array}\right] \ominus\left[\begin{array}{c}
			\mathbf{R}_{w}^{i}\left(\mathbf{p}_{j}^{w}-\mathbf{p}_{i}^{w}+\frac{1}{2} \mathbf{g}\Delta t^{2}-\mathbf{v}_{i}^{w} \Delta t\right) \\
			\mathbf{R}_{w}^{i}\left(\mathbf{v}_{j}^{w}+\mathbf{g} \Delta t-\mathbf{v}_{i}^{w}\right)\\
			\mathbf{q}_{i}^{w_{}^{-1}} \otimes \mathbf{q}_{j}^{w_{}}\\
			\mathbf{b}_{a {i}}\\
			\mathbf{b}_{w {i}}\\
		\end{array}\right]
	\end{aligned}
\end{equation}

Where rotation matrices $\mathbf{R}$ and Hamilton quaternions $\mathbf{q} $ are used to represent rotation. $\mathbf{p}_{i}^{w}$, $\mathbf{v}_{i}^{w}$ and $\mathbf{q}_{i}^{w}$ are translation, velocity, and rotation from the body frame to the world frame. $\mathbf{g}$ is the gravity vector in the world frame.
$\mathbf{b}_{a}$ and $\mathbf{b}_{w}$ are accelerometer bias and gyroscope bias respectively. $\ominus$ is the minus operation on the IMU residuals.  
% The first three rows represent the relative position error, velocity error and rotation error between two nodes. The last two rows represent the accelerometer
Note that the accelerometer bias error and gyroscope bias error are jointly optimized in the graph.   

\subsubsection{\textbf{Relative Pose Factor}}
Since the camera and LiDAR are not globally referenced and their odometry is based on the first pose of the robot, we only use their relative poses [$\Delta p_j^i$, $\Delta q_j^i$] as local constraints to constrain the IMU preintegration's position $\hat \alpha _j^i$, and rotation $\hat \gamma _j^i$. 
The relative pose factor $\mathbf{e}_{i j}^{LIO},\mathbf{e}_{i j}^{VIO}$ is derived as:

\begin{equation}
	\mathbf{e}_{i j}^{LIO},\mathbf{e}_{i j}^{VIO}=\left[ {\begin{array}{*{20}{c}}
{\Delta p_j^i}\\
{\Delta q_j^i}
\end{array}} \right] \ominus \left[ {\begin{array}{*{20}{c}}
{\hat \alpha _j^i}\\
{\hat \gamma _j^i}
\end{array}} \right]
\end{equation}

\subsubsection{\textbf{IMU Odometry Optimization}}
For each new frame, we minimize an optimization problem that consists of relative pose factor ($\mathbf{e}_{i j}^{LIO}$,$\mathbf{e}_{i j}^{VIO}$), and marginalization prior $\mathbf{E}_{m}$.
\begin{equation}
	E=\sum_{(i, j) \in \mathcal{B}} \mathbf{e}_{i j}^{LIO\top} {W}_{i j}^{-1} \mathbf{e}_{i j}^{LIO}+
	\sum_{(i, j) \in \mathcal{B}} \mathbf{e}_{i j}^{VIO\top}{W}_{i j}^{-1} \mathbf{e}_{i j}^{VIO}+\mathbf{E}_{\mathrm{m}}
\end{equation}
The relative pose factor has to be weighted with the appropriate covariance matrix ${W}_{ij}$, which can be calculated based on the reliability of the observation. For example, in visually degraded environments, $\mathbf{e}_{ij}^{VIO}$ calculated by VIO will have a lower weight. In contrast, in geometrically degraded environments, $\mathbf{e}_{ij}^{LIO}$ calculated by LIO will have a lower weight. Note that without loss of generality, the IMU odometry can also incorporate measurements from other sensors such as GPS and wheel odometry. 

% As long as the bias could be constrained by other sensors, Super Odometry will achieve robust state estimation.


\subsection{LiDAR-Inertial Odometry Factors}
The structure of LiDAR-Inertial odometry factors is shown in Fig.\ref{fig:algorithm}.({c}). Since IMU odometry is computationally efficient and outputs at a very high frequency, its IMU preintegration factor can be naturally added into the factor graph of LiDAR-Inertial odometry. The IMU preintegration factor will be used as motion prediction for current scan-map matching and connects consecutive LiDAR frames in the factor graph. 

% \todo{need to add sub pipeline figure ?}

The pipeline of the scan-map matching process can be divided into four steps, namely: (a) PCA-based Feature Extraction (b) Multi-metric Linear Square ICP (c) LiDAR-inertial Odometry Factor (d) Dynamic Octree.

\subsubsection{\textbf{PCA Based Feature Extraction}}
When the system receives a new LiDAR scan, the LiDAR scan will be downsampled to a fixed number and then fed into the feature extraction module. The point-wise K-D tree is used to find the nearest $k$ neighbor points within a sphere of radius $r$. Based on these neighbor points, the principal components analysis (PCA) method is used to analyze the local linearity $\sigma_{1 \mathrm{D}}$, planarity $\sigma_{2 \mathrm{D}}$, and curvatures $\sigma_{3 \mathrm{D}}$ of geometrical feature. They are defined as\cite{hackel2016fast}
\begin{equation}
	\sigma_{1 \mathrm{D}}=\frac{\sigma_{1}-\sigma_{2}}{\sigma_{1}}, \sigma_{2 \mathrm{D}}=\frac{\sigma_{2}-\sigma_{3}}{\sigma_{1}}, \sigma_{3\mathrm{D}}=\frac{\sigma_{3}}{\sigma_{1}+\sigma_{2}+\sigma_{3}}
	\label{eq:feature}
\end{equation}
where $\sigma_{i}=\sqrt{\lambda_{i}}$ and $\lambda_{i}$ are the eigenvalues of the PCA for the normal computation. According to the magnitude of local feature $\sigma_{1 \mathrm{D}}$,$\sigma_{2 \mathrm{D}}$, $\sigma_{3 \mathrm{D}}$, point, line and plane features $\mathbb{F}_{i+1}=\{F^{po}_{i+1}, F^{li}_{i+1},F^{pl}_{i+1}\}$ can be easily classified.

\subsubsection{\textbf{Multi-metric ICP Factors}}
After extracting reliable geometric features,
we use motion prediction provided by IMU odometry and transform the $\mathbb{F}_{i+1}$ from $\bold{B}$ to $\bold{W}$. Then, we find its point, line, and plane correspondence in the map by using the proposed dynamic octree. For the sake of brevity, the detailed procedures will be described in the next section.

To improve the robustness of Super Odometry, we expect to estimate the optimal transformation that jointly minimizes the point-to-point, point-to-line, and point-to-plane distance error metric, which can be computed using the following equations:
\begin{equation}
	\begin{cases}
		\mathbf{e}_{i}^{\mathrm{po} \rightarrow \mathrm{po}}=\mathbf{q}_{i}-(R\mathbf{p}_{i}+t) & \text{$q_{i},p_{i}\in F^{po}_{i+1}$}
		\\
		\mathbf{e}_{j}^{\mathrm{po} \rightarrow \mathrm{li}}=\mathbf{v}_{j} \times\left(\mathbf{q}_{j}-(R\mathbf{p}_{j}+t)\right)
		& \text{$q_{j},p_{j}\in F^{li}_{i+1}$}
		\\
		\mathbf{e}_{k}^{\mathrm{po} \rightarrow \mathrm{pl}}=\mathbf{n}_{k}^{\top}\left(\mathbf{q}_{k}- (R\mathbf{p}_{k}+t)\right) 
		& \text{$q_{k},p_{k}\in F^{pl}_{i+1}$}
	\end{cases}
\end{equation}

where $\mathbf{e}_{i}^{\mathrm{po} \rightarrow \mathrm{po}}$, $	\mathbf{e}_{j}^{\mathrm{po} \rightarrow \mathrm{li}}$ and $\mathbf{e}_{k}^{\mathrm{po} \rightarrow \mathrm{pl}}$ are the point-to-point (line, plane) distance. $\mathbf{v}_{j}$ and $\mathbf{n}_{k}$ represent the corresponding principal and normal direction of features. $\mathbf{T}_{i+1}=\{\bold{R}, \bold{t}\}$ is the optimal transformation we expect to estimate. 

However, in airborne obscurants environments such as dust, fog, and smoke environments, we may extract geometric features on these noise data and obtain unreliable data association. To solve this problem, we found that the quality of correspondence is depended on if the extracted features' neighbor points fit the point, line, or plane distribution. Therefore, we use distribution's quality ${W}_{l}=\{w_i^{po \to po}, w_i^{po \to li},w_i^{po \to pl}\}$ to define the quality of correspondence by using the following equation.  

\begin{equation}
\label{eq:w1}
{w_i^{po \to po} = \frac{{{\sigma _{3D}}}}{{{\sigma _{3Dmax}}}}}
\end{equation}


\begin{equation}
\label{eq:w2}
{w_i^{po \to li/pl} = 1 - \frac{{\sum\limits_{i = 1}^k {{{\left( {{p_i} - \bar p} \right)}^T}A\left( {{p_i} - \bar p} \right)} }}{{k \times {d_{max}}}}}
\end{equation}

where $\sigma_{3Dmax}$ is the maximum threshold of curvatures.
$p_{i}$ are the neighbor points of extracted features and $\bar p$ is their mean. For point-to-line(plane) distance error, we use normal direction $\mathbf{n}_{k}$ of features and define $A=I-\mathbf{n}_{k}^T \times \mathbf{n}_{k}$ and $A=\mathbf{n}_{k}^T \times \mathbf{n}_{k}$ respectively.
$d_{max}$ represents the maximum distance error for point-to-line(plane) correspondence. $k$ is the number of neighbor points.  

\subsubsection{\textbf{LiDAR-inertial Odometry Optimization}} The levenberg-marquardt method is then used to solve for the optimal transformation by minimizing:

\begin{equation}
\mathop {\min }\limits_{{\mathbf{T}_{i + 1}}} \left\{ {\left. {\begin{array}{*{20}{c}}
{\sum\limits_{p \in {\mathbb{F}_{i}}} {{{\left\| {{W_l}e_i^{po \to po,li,pl}} \right\|}^2} + } }\\
{\sum\limits_{\left( {i,i + 1} \right) \in \mathcal{B} } {{{\left\| {{W_{imu}}{e^{imu}}} \right\|}^2} + \mathbf{E}_{imuodom}^{prior}} }
\end{array}} \right\}} \right.
\end{equation}

$\mathbf{E}_{imuodom}^{prior}$  and $\mathbf{e}_{}^{imu}$ are the predicted pose prior and IMU preintegration factor provided by IMU odometry. $\mathbf{W}_{l}$ and $\mathbf{W}_{imu}$ are the weight of LiDAR and imu correspondence.  

When the environment is geometrically degraded, since the IMU odometry can receive the measurements from other sensors, the IMU odometry is still reliable and $E_{imuodom}^{prior}$ will be dominant in the optimization problem. Meanwhile, unreliable LiDAR constraints will be rejected or have low weight by evaluating the quality of constraints. On the other hand, when the environment is well-structured, the multi-metric ICP factor will be dominant in the optimization problem. 


% \subsubsection{\textbf{Dynamic Octree}}
% Most LiDAR SLAM methods adopt the KD-tree method to organize the 3D points and fulfill data association. However, we argue that the traditional KD-tree is very time-consuming to organize the 3D points since it only uses one tree to organize all the points and needs to recreate the KD-tree every time when adding new points shown in Fig.\ref{fig:dynamic_octree}(b). 
% % \textbf{Key Insight}: Instead of improving the neighbor search in the three-dimensional point sets, we push more effort on how to reuse the structure of the existing tree and avoid recreating a new tree when the map is updated. To achieve this, we propose a more efficient organization of 3D points named dynamic octree to facilitate data association.  
% \begin{figure}
% 	\centering
% 	\includegraphics[width=0.85\columnwidth]{Method/dynamic_octree.jpg}
% 	\caption{The comparison of 3D KD-tree and dynamic octree. (a,b) shows the construction process of 3D KD-tree when adding new point cloud (grey circle), which needs to change the structure of the whole tree. (c,d) shows the construction process of dynamic octree when adding new point cloud (grey circle), which only needs to change the structure of the sub-tree.}
% 	\vspace{-3mm}
% 	\label{fig:dynamic_octree}
% \end{figure}
%   To solve this problem, we propose to use  a more efficient organization of 3D points named dynamic octree to facilitate data association. The dynamic octree is based on the work of Behley \cite{behley2015icra}. It stores the map as a hash table as shown in Fig.\ref{fig:dynamic_octree}(d). 
% The world is represented by voxels and voxels can be accessed from the hash table by using their XYZ indices. Instead of only building one tree, each voxel will have its octree to organize the points and each octree can be accessed by a hash table. Since we only need to update the specific octree instead of the whole tree, the real-time performance on data association will be significantly improved.

\subsection{Visual-Inertial Odometry Factors}
To take full advantage of the fusion of vision and LiDAR sensing modalities, we track monocular visual features and LiDAR points within the camera field of view and use them to provide depth information for visual features.


\subsubsection{ \textbf{Visual-inertial Odometry Optimization}}
For each new keyframe, we minimize a non-linear optimization problem that consists of visual reprojection factor $\mathbf{e}_{reproj}$, IMU preintegration factor $\mathbf{e}_{imu}$, marginalization factor $\mathbf{E}_{m}$ and a pose prior from IMU odometry $\mathbf{E}_{imuodom}^{prior}$. The structure of Visual-inertial odometry factors is shown in Fig.\ref{fig:algorithm}(b).

\begin{equation}
	\min _{\mathbf{T}_{i+1}}\left\{
	\begin{split}
		&\sum_{ i \in \mathbf{obs}(i)}\mathbf{e}_{\text {reproj }}^{\mathbf{T}} W_{\text {reproj}}^{-1} \mathbf{e}_{\text {reproj }}^{\mathbf{}}+\sum_{(a, b) \in \mathbf{C}}\mathbf{e}_{\text {imu}}^{\mathbf{T}} W_{imu}^{-1} \mathbf{e}_{imu}^{\mathbf{}}
		\\ 
		&+\mathbf{E}_{m} + \mathbf{E}_{imuodom}^{prior}	
	\end{split}
	\right\}
\end{equation}

$\mathbf{obs(i)}$ represents a set that contains visual features $i$ that is tracked by other frames. The set $\mathbf{C}$ contains pairs of visual nodes $(a, b)$ connected by IMU factors. $W_{\text {reproj}}$ and $W_{\text {imu}} $ are covariance matrices for visual reprojection factor and IMU preintegration factor. 

Similar to the explanation in the last section, when the environment is visually degraded, the $\mathbf{E}_{imuodom}^{prior}$ will be dominant in the optimization problem and unreliable vision factor will be rejected by analyzing the information matrix. However, when the environment has good lighting conditions, the visual factor will be dominant in the optimization problem and the IMU preintegration factor only provides the initial guess for visual feature tracking.   